% ==========================================
% SEZIONI DEL VERBALE
% ==========================================
\newpage
\vspace{0.8cm}

\section{Ordine del giorno}

\subsection{Presentazione e approccio al progetto}
Durante l'incontro, il team ha illustrato l'impostazione logica preliminare e la modellazione teorica relative alla pagina di login e all'architettura multi-tenant. Il referente aziendale ha confermato l'approvazione del template concettuale di autenticazione, ribadendo che in questa fase di studio e definizione dei requisiti l'attenzione deve rimanere focalizzata esclusivamente sulla logica di funzionamento e sui flussi funzionali, tralasciando i dettagli legati all'User Experience.
\subsection{Gestione degli store Shopify e primo use case (RF001)}
In merito ai flussi operativi successivi all'autenticazione, il referente ha chiarito le modalità di interazione con la piattaforma e-commerce:
\begin{itemize}
    \item \textbf{Funzionalità di registrazione:} Il primo caso d'uso operativo da implementare prevede la presenza di un apposito pulsante denominato ``Registra store''.
    \item \textbf{Flusso di configurazione:} L'utente inserisce l'URL dello store, procedendo successivamente alla generazione del token e al salvataggio sicuro delle chiavi API di Shopify.
    \item \textbf{Modello dati e scalabilità:} Per la fase iniziale è previsto il supporto a un singolo store per ciascuna organizzazione; tuttavia, il modello dei dati sottostante deve essere opportunamente predisposto per supportare la gestione di store multipli in futuro.
\end{itemize}
\subsection{Roadmap dei Requisiti}
Durante la discussione sulla pianificazione evolutiva delle funzionalità e sulla definizione dei casi d'uso derivati dalla dashboard, sono state approfondite le aspettative relative alla visualizzazione e alle interazioni utente:
\begin{itemize}
    \item \textbf{Scomposizione di RF005:} È stato stabilito di suddividere il requisito RF005 in sottorequisiti puntuali e specifici. A titolo esemplificativo, vengono definiti:
    \begin{itemize}
        \item \texttt{RF005.1} per l'analisi del catalogo e dei metafield;
        \item \texttt{RF005.2} per la gestione dei webhook;
        \item \texttt{RF005.3} per il controllo delle app installate e del tema/storefront.
    \end{itemize}
    \item \textbf{Sequenza del percorso logico:} La roadmap consigliata per la verifica end-to-end prevede una progressione lineare che conduce da: Login $\rightarrow$ Tenant $\rightarrow$ Store $\rightarrow$ Scansione di catalogo e metafield $\rightarrow$ Analisi dell'utilizzo dei metafield mediante il requisito \texttt{RF007}.
    \item \textbf{Derivazione dei casi d'uso per la dashboard:} L'analisi delle interfacce e dei flussi operativi attesi sulla dashboard consente di strutturare i casi d'uso focalizzandosi sulle funzionalità core e sulla visualizzazione immediata dello stato dello store, privilegiando la concretezza delle informazioni rispetto a elementi superflui di interfaccia.
\end{itemize}
\subsection{Architettura Multi-Tenancy e gestione del database PostgreSQL}
È stata esaminata l'implementazione del multi-tenancy e dell'isolamento dei dati:
\begin{itemize}
    \item \textbf{Isolamento e Tenant:} Ciascun utente deve operare all'interno di un ambiente isolato (organizzazione o tenant). Non è richiesta una complessa gerarchia di permessi intra-tenant nella fase iniziale, in quanto tutti gli utenti appartenenti alla stessa organizzazione condividono gli stessi privilegi.
    \item \textbf{Row Level Security (RLS):} È stata confermata l'impostazione di base tramite PostgreSQL e l'utilizzo di wrapper di transazioni per abilitare la sicurezza a livello di riga (RLS), sebbene il concetto di organizzazione debba essere pienamente integrato nello schema del database.
    \item \textbf{Modello Dati:} Si è convenuto di strutturare il database introducendo una tabella dedicata alle organizzazioni in relazione uno-a-molti con gli utenti, evitando di utilizzare il dominio email come chiave del tenant. Per la gestione e il popolamento iniziale, è stato suggerito l'utilizzo di script tramite Drizzle ORM.
\end{itemize}
\subsection{Gestione di Vercel e Flusso di Deploy CI/CD}
In merito all'infrastruttura di deployment, il referente ha confermato che il team ha attivato un account a pagamento sul piano Vercel Pro finanziato dall'azienda tramite una mail di gruppo:
\begin{itemize}
    \item \textbf{Workflow:} Il flusso operativo prevede l'associazione di Vercel direttamente a GitHub tramite le GitHub Actions. Per ogni push su un branch di feature, Vercel genera automaticamente una preview di build, eliminando la necessità di accedere manualmente alla piattaforma. Sulle Pull Request verso il branch principale viene eseguita la build di produzione.
    \item \textbf{Account condiviso:} Considerato il costo degli account Vercel, l'utilizzo di un'unica email condivisa per il team rappresenta la soluzione ottimale per la gestione della pipeline.
\end{itemize}
\subsection{Calendario prossimi incontri}
È stato definito un primo calendario di incontri che si svolgeranno a cadenza settimanale di mercoledì su Google Meet dalle ore 14 alle ore 15:
\begin{itemize}
    \item \textbf{25 Novembre 2026:} \textit{Incontro saltato} (azienda in evento aziendale offsite).
    \item \textbf{02 Dicembre 2026 -- 09 Dicembre 2026:} Incontri ordinari di avanzamento.
    \item \textbf{15 Dicembre 2026:} Anticipo rispetto a mercoledì 16 per disponibilità referente.
    \item \textbf{07 Gennaio 2027 -- 13 Gennaio 2027:} Incontri finali pre-consegna POC (prevista per metà Gennaio).
\end{itemize}

